   %% ============================================================
%%  Section 3 — Problem Definition
%%  Article: BRKGA with Rolling Horizon for the MDRP
%%  Citation style: LNCS (numeric \cite{})
%% ============================================================

\section{Problem Definition}\label{sec:problem}

We study a single-platform, multi-courier meal delivery operation
in a dense urban environment.
Each day the platform receives a stream of customer orders that
must be collected from restaurants and delivered to customers,
while respecting courier working schedules, vehicle capacities,
and time-window commitments.
Following the formulation of Reyes et al.~\cite{reyes2018mdrp},
we model this as an online combinatorial optimisation problem
in which decisions are made sequentially as new information
becomes available.

%% ------------------------------------------------------------------
\subsection{Sets and Parameters}\label{subsec:notation}
%% ------------------------------------------------------------------

Table~\ref{tab:notation} summarises the sets and parameters used
throughout the paper.

\begin{table}[ht]
\centering
\caption{Sets and parameters of the MDRP.}
\label{tab:notation}
\small
\setlength{\tabcolsep}{5pt}
\begin{tabular}{lp{9.5cm}}
\toprule
\textbf{Symbol} & \textbf{Definition} \\
\midrule
\multicolumn{2}{l}{\textit{Sets}} \\
$U$   & Set of customers; each $u \in U$ has delivery location $\ell_u$ \\
$R$   & Set of restaurants; each $r \in R$ has location $\ell_r$ \\
$O$   & Set of orders; each $o \in O$ is associated with restaurant $r_o \in R$,
        customer $u_o \in U$, preparation time $e_o$, ready time $s^r_o$,
        and deadline $s^u_o$ \\
$C$   & Set of couriers; each $c \in C$ has location $\ell_c$,
        log-in time $\alpha_c$, and log-out time $\delta_c$ \\
$S$   & Set of feasible routes; each $s \in S$ has one pickup node
        $r_s \in R$ and an ordered sequence of delivery nodes
        $U_s \subseteq U$ \\
$T$   & Set of vehicle types
        (motorcycle, bicycle, car, pedestrian) \\
\midrule
\multicolumn{2}{l}{\textit{Parameters}} \\
$v_c$             & Speed of courier $c$ (km/min), determined by vehicle type \\
$d_{ij}$          & Haversine distance between locations $i$ and $j$ (km) \\
$t_{ij}$          & Travel time between $i$ and $j$: $t_{ij} = d_{ij}/v_c$ \\
$p_o$             & Compensation paid to the courier for delivering order $o$ \\
$K$               & Maximum number of orders per route \\
$D_{\max}$        & Maximum distance from a courier to a pickup location for
                    eligibility (km) \\
$\Delta$          & Rolling-horizon window length (min) \\
\bottomrule
\end{tabular}
\end{table}

%% ------------------------------------------------------------------
\subsection{Decision Structure and Rolling Horizon}\label{subsec:decisions}
%% ------------------------------------------------------------------

The planning day is divided into consecutive, non-overlapping
windows of length $\Delta$ minutes.
At the beginning of each window $w$, the platform observes:

\begin{itemize}
    \item the set of \emph{pending orders} $O^w \subseteq O$—orders
    that are ready or will be ready within the current window and
    have not yet been assigned;
    \item the set of \emph{available couriers} $C^w \subseteq C$—
    couriers that are idle or completing their current route before
    the window ends.
\end{itemize}

\noindent
The platform then selects a set of \emph{assignments}
$\mathcal{A}^w = \{(c, s) : c \in C^w,\; s \in S^w\}$,
where $S^w$ is the set of routes constructed from $O^w$,
subject to the constraints detailed in
Section~\ref{subsec:constraints}.
Each assignment is communicated to the corresponding courier
at the start of window $w$; no reassignment is possible once
a courier departs for the pickup location.
The plan is then \emph{rolled forward}: unassigned orders
and couriers whose routes finish before the next window
re-enter the eligible pools at window $w+1$.

This structure decouples the global planning problem into a
sequence of short-horizon sub-problems, each of which is
solved by one of the four methods compared in this paper
(Section~\ref{sec:methodology}).

%% ------------------------------------------------------------------
\subsection{Operational Constraints}\label{subsec:constraints}
%% ------------------------------------------------------------------

The following constraints apply uniformly to all three methods:

\begin{enumerate}
    \item \textbf{Route capacity.}
    Each route $s$ contains at most $K$ orders, all originating
    from the same restaurant: $|U_s| \leq K$.

    \item \textbf{Courier eligibility.}
    A courier $c$ is eligible for route $s$ only if
    $d_{\ell_c, \ell_{r_s}} \leq D_{\max}$ and the courier's
    log-out time satisfies $\delta_c \geq s^r_{o}$
    for all $o \in U_s$.

    \item \textbf{Time-window feasibility.}
    For a route assigned to courier $c$, the courier must arrive
    at each customer node $u_i \in U_s$ no later than the
    deadline $s^u_{o_i}$.
    Arrival times are propagated sequentially:
    \begin{equation}
        a_{u_1} = \alpha_{r_s} + t_{\ell_{r_s}, \ell_{u_1}},
        \qquad
        a_{u_{i+1}} = a_{u_i} + s_{u_i} + t_{\ell_{u_i}, \ell_{u_{i+1}}},
    \end{equation}
    where $\alpha_{r_s}$ is the courier's arrival time at the
    pickup restaurant and $s_{u_i}$ is the service time at node $u_i$.

    \item \textbf{Single assignment.}
    Each order is assigned to at most one courier,
    and each courier executes at most one route per window:
    \begin{equation}
        \sum_{s \in S^w} \mathbf{1}[(c,s) \in \mathcal{A}^w] \leq 1
        \quad \forall c \in C^w.
    \end{equation}

    \item \textbf{Notification rules.}
    A courier can only be notified of a new assignment when
    idle or en route to a pickup location.
    Couriers in transit to a customer node are not interrupted.

    \item \textbf{Cancellations.}
    Once an order is cancelled by the platform or the customer,
    it is permanently removed from $O$ and cannot be reassigned.
\end{enumerate}

%% ------------------------------------------------------------------
\subsection{Performance Metric}\label{subsec:objectives}
%% ------------------------------------------------------------------

The primary performance metric throughout this paper is the
\emph{coverage rate}: let $\bar{O} \subseteq O$ denote the set
of orders successfully delivered by the end of the planning day;
then

\begin{equation}\label{eq:coverage}
    \rho = \frac{|\bar{O}|}{|O|}.
\end{equation}

\noindent
A higher $\rho$ reflects greater platform reliability and
customer satisfaction.
The four methods evaluated in Section~\ref{sec:results} all
target the maximisation of $\rho$, subject to the operational
constraints of Section~\ref{subsec:constraints}.

\paragraph{Remark on coverage as an objective rather than a constraint.}
One might alternatively impose a service-level constraint of the
form $\rho \geq \rho^{\min}$ (e.g., $\rho^{\min} = 0.95$).
We treat coverage as an \emph{objective} for two reasons.
First, in a rolling-horizon setting the platform cannot know at
dispatch time whether a target end-of-day coverage is achievable:
a constraint that is satisfiable over the full horizon may be
locally infeasible at every intermediate window, preventing any
assignment from being dispatched.
Second, the real instances used here contain high-volume days on
which even the best methods reach $\rho \approx 34$--$47\%$,
reflecting structural fleet-capacity limitations rather than
algorithmic shortcomings; imposing an SLA above this ceiling
would render the sub-problem infeasible.
Treating coverage as an objective is consistent with the
established MDRP literature~\cite{reyes2018mdrp,ulmer2021restaurant}
and always yields a feasible dispatch.
